\section{Sketch Propagation Framework}
\label{sec:scopeprop}

In this section, we introduce a novel \emph{sketch propagation} framework to efficiently estimate the degree of each node in $H$ by constructing cardinality sketches for their neighbors. As a result, it effectively answers \ced queries based on degree centrality.
In addition, we propose an early termination technique to further accelerate personalized \ced query processing.

\subsection{Sketch Propagation for \ced}
\label{sec:sketch-framework}

Our basic idea is to approximate the degree of each node in $H$ by constructing KMV sketches for the neighborhood $\nset(v)$ of each node $v\in V$ without materializing $H$. The KMV sketch was initially proposed for distinct-value estimation \cite{beyer2007synopses} and defined as follows.
\begin{definition}[KMV Sketch]\label{def:kmv}
  Given a collection $C=\{o_0, o_1,$ $ \dots, o_{|C|}\}$ of items and an integer $k > 0$, the KMV ($k$-th Minimum Value) sketch $\kset(C)$ is built by independently drawing a number uniformly at random from $(0, 1)$ for each item in $C$ and maintaining the $k$ smallest random numbers. The set of random numbers generated for each item in $C$ is called the \emph{basis} for the KMV sketch. The cardinality of $C$ is estimated as $|\widetilde{C}|=\frac{k}{\mathbb{E}[\zeta]}-1$, where $\zeta$ is a random variable by taking the largest number in $\kset(C)$.
\end{definition}

By constructing multiple (specifically $\theta$) KMV sketches and applying the Bernstein inequality~\cite{audibert2007tuning}, an estimation of the collection size with high probability can be obtained. 
\begin{lemma}\label{theorem:dconcentration}
  For any $\delta, p \in (0, 1)$, if $\theta = \Theta(\frac{|C|}{\delta k}\log{\frac{1}{p}})$, we have $(1 - \delta) |C| \leq |\widetilde{C}| \leq (1+\delta) |C|$ with probability at least $1 - p$.
\end{lemma}

Note that the KMV sketch is designed primarily for stream processing \cite{estan2003bitmap, alon1996space, durand2003loglog, heule2013hyperloglog, gibbons2001distinct}, where a one-pass scan over the collection $C$ is required.
However, for \ced queries, scanning the neighborhood of each node is quite expensive, as discussed in the description of exact algorithms.
To address this challenge, we propose the \emph{sketch propagation} framework that constructs a KMV sketch for $\nset(v)$ of each $v \in V$ without explicitly generating and scanning $\nset(v)$ and offers a different estimator with a tight error bound for the case where $\nset(v)$ is relatively small based on the notion of \emph{images}.
\begin{definition}[Images]
  The image of a matching node $u \in \vset_i$, denoted as $\iset(u)$, contains a node $v \in \vset_0$ if there exists an instance $M$ of $\mset$ including both $u$ and $v$, i.e., 
  $$\iset(u) = \{ v \in \vset_0 \mid \exists M \triangleright \mset, M(x_0) = v \wedge M(x_i) = u \}.$$
\end{definition}
\begin{example}
  In Fig.~\ref{fig:propagation}$(a)$, the image $\iset(p_1)$ contains two nodes $a_2$ and $a_3$, as they are connected with $p_1$ via $M(a_2, p_1, v_0, \bar{p}_0,$ $\bar{a}_0)$ and $M(a_3, p_1, v_0, \bar{p}_0, \bar{a}_0)$, respectively.
\end{example}
An important observation is that $\nset(u) = \iset(\bar{u})$ for each node $u \in V=\vset_0$ in the relational graph due to its definition.

\vspace{1mm}
\myp{Sketch Propagation}
We build a KMV sketch for $\nset(u)$ of each $u \in V$ by iteratively maintaining the sketches for the images of nodes from $\vset_0$ to $\vset_L$.
First, a random number $r(u)$ is generated for each node $u \in \vset_0$. Since $\iset(u) = \{u\}$, the sketch is initialized as $\kset(\iset(u))=\{r(u)\}$. 
The sketches for the images of the nodes in $\vset_i$ are then constructed by iteratively propagating the sketches of the nodes in $\vset_{i-1}$.
Specifically, each node in $\vset_{i-1}$ propagates its sketch to all its successor nodes in $\vset_{i}$ and a node in $\vset_i$ builds its sketch by retaining the $k$ smallest random numbers among all sketches it receives.
The following lemma ensures the correctness of the \emph{sketch propagation} framework.

\begin{algorithm}[t]
\caption{Sketch Propagation}\label{alg:propagation}
    \LinesNumbered
    \KwIn{KG $\gset$, meta-path $\mset$, quantile parameter $\lambda$, number of propagations $\theta$, KMV sketch size $k$}
    \KwOut{A set of hubs $S$}
    \ForAll{$u \in \vset^*$}{\label{algline:initK-begin}
        \For{$t=1$ \KwTo $\theta$}{
            $\kset[u][t] \gets \varnothing$\;
            \textbf{if} $u\in\vset_0$ \textbf{do} add $r(u) = \mathtt{Rand}(0, 1)$ to $\kset[u][t]$\;\label{algline:initK-end}
        }
    }
    $\widetilde{N} \gets$ \texttt{Propagate}$(\kset)$\;
    $S \gets$ top-$(\lambda\cdot|V|)$ nodes with highest degrees in $H$ w.r.t.~$\widetilde{N}$\;\label{algline:spreturn-begin}
    \textbf{return} $S$\;\label{algline:spreturn-end}
    
    \BlankLine
    \SetKwFunction{FMain}{Propagate}
    \SetKwProg{Fn}{Function}{:}{}
    \Fn{\FMain{$\kset$}}{
        \For{$i=1$ \KwTo $L$}{ \label{algline:construction-begin}
            \textbf{forall} $u \in \vset_i$ \textbf{do} \texttt{Receive}$(u, \kset)$\;
        }\label{algline:construction-end}
        \ForAll{$u\in\vset_L$}{ \label{algline:estimation-begin}
            $\widetilde{\mu} \gets 0$\;
            \For{$t=1$ \KwTo $\theta$}{
                \uIf{$\left| \kset[u][t] \right| = k$}{
                    $\widetilde{\mu} \gets \widetilde{\mu} + \frac{\max(\kset[u][t])}{\theta}$\;
                }
                \Else{
                    $\widetilde{\mu} \gets \widetilde{\mu} + \frac{k}{(|\kset[u][t]| + 1) \cdot \theta}$\;
                }
            }
            $\widetilde{N}[u] \gets k/\widetilde{\mu} - 1$\;\label{algline:estimation-end}
        }
        \Return{$\widetilde{N}$}\;
    }
    \BlankLine
    \SetKwFunction{FMain}{Receive} 
    \SetKwProg{Fn}{Function}{:}{} 
    \Fn{\FMain{$u, \kset$}}{\label{algline:receive}
        \For{$t=1$ \KwTo $\theta$}{\label{algline:receive-begin}
            $\kset[u][t] \gets \oplus_{v \in \vset^-(u)}{\kset[v][t]}$\;\label{algline:receive-end}
        }
    }
\end{algorithm}

\begin{lemma}\label{lemma:img}
  Given a meta-path $\mset$, for each node $u \in \vset_i$ 
  \begin{equation}\label{equation:img1}
    \kset(\iset(u))=\oplus_{v\in\vset^-(u)}{\kset(\iset(v))} 
  \end{equation}
hold if all KMV sketches are constructed from the same basis on $\vset_0$. The operator $\oplus$ extracts the $k$ smallest random numbers from the union of KMV sketches.
\end{lemma}
\techreport{
\begin{proof}
  According to \cite[Thm.~5]{beyer2007synopses}, $\kset(A) \oplus \kset(B)$ is a KMV sketch for the collection $A \cup B$ for two collections $A$ and $B$.
  Therefore, we only need to prove $\iset(u)=\bigcup_{v\in\vset^-(u)}\iset(v)$ for each $u\in\vset_i$.

  On the one hand, for any node $u^\prime\in\iset(u)$, there exists $M\triangleright\mset$ such that $M(x_i)=u$ and $M(x_0)=u^\prime$. 
  Suppose that $v=M(x_{i-1})$, we have $u^\prime\in\iset(v)$ and $v\in\vset^-(u)$. Thus, $u^\prime\in\bigcup_{v\in\vset^-(u)}\iset(v)$, i.e., $\iset(u)\subseteq\bigcup_{v\in\vset^-(u)}\iset(v)$. 

  On the other hand, for any node $u^\prime\in \bigcup_{v\in\vset^-(u)}\iset(v)$, there exists a node $v\in\vset^-(u)$ such that $u^\prime\in\iset(v)$, i.e., there exists an instance $M^\prime\triangleright\mset$ such that $M^\prime(x_0)=u^\prime$ and $M^\prime(x_{i-1})=v$.
  Moreover, since $v\in\vset^-(u)$, there exists an instance $M^{\prime\prime}\triangleright\mset$ such that $M^{\prime\prime}(x_{i-1})=v$ and $M^{\prime\prime}(x_i)=u$. 
  By concatenating $M^\prime(x_0)=u^\prime$ to $M^\prime(x_{i-1})=v$ with $M^{\prime\prime}(x_{i-1})=v$ to $M^{\prime\prime}(x_L)$, we construct an instance $M\triangleright\mset$ such that $M(x_0)=u^\prime$ and $M(x_i)=u$. Thus, $u^\prime\in\iset(u)$, i.e., $\bigcup_{v\in\vset^-(u)}\iset(v)\subseteq\iset(u)$.
  
  Thus, we have $\iset(u)=\bigcup_{v\in\vset^-(u)}\iset(v)$ for each $u\in\vset_i$.
\end{proof}
}

Algorithm~\ref{alg:propagation} details the \emph{sketch propagation} framework.
It receives a KG $\gset$, a meta-path $\mset$, a parameter $\lambda \in (0, 1)$, the number of propagations $\theta \in \mathbb{Z}^+$, and the sketch size $k \in \mathbb{Z}^+$ as input, and returns a set of nodes $S$ for the (approximate) \ced query.
It first initializes an array $\kset$ to store $\theta$ KMV sketches for images of each node $u \in \vset^*$ (Lines~\ref{algline:initK-begin}--\ref{algline:initK-end}).
Next, it calls the function \texttt{Propagate}, which constructs KMV sketches within $\kset$ (Lines~\ref{algline:construction-begin}--\ref{algline:construction-end}) and estimates the degree of each node $u$ in $H$ using $\kset$ (Lines~\ref{algline:estimation-begin}--\ref{algline:estimation-end}).
Specifically, the function \texttt{Receive} depicts the step for KMV sketch construction when each node receives the sketches propagated from other nodes.
Given a node $u$, the array $\kset$ to store sketches, it scans the random numbers in the sketches of all nodes in $\vset^-(u)$ and keeps the $k$ smallest numbers in $\kset[u][t]$ with a min-heap of size $k$ (Lines~\ref{algline:receive-begin} and~\ref{algline:receive-end}).
Finally, it returns top-$(\lambda \cdot |V|)$ nodes with the highest estimated degrees as $S$ (Lines~\ref{algline:spreturn-begin} \&~\ref{algline:spreturn-end}). The ties are broken arbitrarily.

\begin{figure}[t]
    \centering
    \includegraphics[width=8cm]{figure/forward.pdf}
    \vspace{-1em}
    \caption{Running example of sketch propagation framework with $k=2$ and $\theta=1$. Subfigure~$(a)$ illustrates the complete sketch propagation process for a \ced query; Subfigure~$(b)$ demonstrates the early terminated sketch propagation for a personalized \ced query with $a_4$ as the query node. The redlined elements are discarded due to the $k$-size limitation.}
    \label{fig:propagation}
\end{figure}

\begin{example}
\label{example:propag}
  Fig.~\ref{fig:propagation}$(a)$ illustrates the sketch propagation process with $k=2$ and $\theta=1$. To begin with, we generate a random number for each node in $\vset_0$ and then propagate random numbers along the matching graph to $\vset_4$. The sketch $\kset(\iset(v_0))$ is constructed by taking the two smallest numbers from the union of sketches of the images of the nodes in $\vset^-(v_0)$. Thus, $\kset(\iset(v_0))=\{0.4, 0.9\}\oplus\{0.2, 0.6\}=\{0.2, 0.4\}$ and the cardinality of $\iset(v_0)$ is estimated as $\widetilde{I}(v_0)=\frac{2}{0.4}-1=4$.
\end{example}

\myp{Theoretical Analysis}
Next, we will prove that \emph{sketch propagation} approximates \ced queries with high probability (w.h.p.) (cf.~Theorem~\ref{theorem:dglobal}).
Hereafter, we use $I(u)$ to denote the sizes of $\iset(u)$ for any node $u\in\vset^*$.
We also use $\widetilde{I}(u)$ to denote the estimate of $I(u)$ provided by sketch propagation.

\techreport{
Moreover, we use $\zeta$ to denote the random variable that takes the maximum random number in the sketch of either $u$'s forward or backward image and $\mu$ to denote the expectation of $\zeta$. 

\begin{lemma}\label{theorem:estimator}
  For any $u\in\vset^*$, $I = k/\mu - 1$ if $I \geq k$.
\end{lemma}
\begin{proof}
According to \cite{beyer2007synopses}, if $I\geq k$, the probability density function of $\zeta$ will be $f_\zeta(t) = \frac{t^{k-1}(1-t)^{I-k}}{B(k, I-k+1)}$,
where $B(a,b)$ denotes the beta function.
Thus, we have
\begin{equation}
\begin{split}
\mu & = \int_{0}^{1} t f_\zeta(t) \,dt = \int_{0}^{1} \frac{t^k(1-t)^{I-k}}{B(k, I-k+1)}\,dt \\
    & = \frac{B(k+1, I-k+1)}{B(k, I-k+1)} = \frac{k\cdot\tbinom{I}{k}}{(k+1)\cdot\tbinom{I+1}{k+1}} = \frac{k}{I+1},
\end{split}
\end{equation}
which indicates that $I = k/\mu - 1$.
\end{proof}

\myp{Remark}
When $I < k$, a KMV sketch directly provides the true value of $I$ with the number of random numbers in the sketch. Furthermore, the original estimator $\mathbb{E}[(k-1)/\zeta]$ in Definition~\ref{def:kmv} is unbiased, but with only asymptotic error bounds. In contrast, our proposed estimator in Lemma~\ref{theorem:estimator} utilizes Bernstein's inequality to establish a tighter confidence bound for arbitrary image sizes.

\begin{theorem}[Bernstein Inequality \cite{audibert2007tuning}]
  Let $X_1$, $X_2$, \dots, $X_\theta$ be real-valued i.i.d.~random variables such that $X_i\in[0, r]$, $\mathbb{E}[X_i]=\mu$, and $\mathrm{Var}[X_i]=\sigma^2$. Let $\overline{X}_\theta=\frac{1}{\theta}\sum^{\theta}_{i=1}{X_i}$ denote the empirical mean. With probability at least $1-p$, we have
  $$|\overline{X}-\mu| \leq \sqrt{\frac{\sigma^2\log(2/p)}{\theta}}+\frac{2r\log(1/p)}{\theta}.$$
\end{theorem}
}

\techreport{
\begin{proof}
When $I < k$, the exact value of $I$ will always be obtained, i.e., $\widetilde{I} = I$, and the lemma holds trivially.
When $I \geq k$, let $\Delta\mu = |\widetilde{\mu}-\mu|$ and $\Delta I = |\widetilde{I}-I|$. According to Lemma~\ref{theorem:estimator}, $I=\frac{k}{\mu} - 1$ and thus 
\begin{equation*}
\Delta I =
\begin{cases}
\frac{k}{\mu} - \frac{k}{\mu + \Delta\mu}, & \text{if } \mu < \widetilde{\mu}; \\
\frac{k}{\mu - \Delta\mu} - \frac{k}{\mu}, & \text{if } \mu \geq \widetilde{\mu}.
\end{cases}
\end{equation*}
For any $\delta \in (0,1)$, it follows that $\Delta I \leq \delta I$ if $\Delta\mu \leq \frac{\delta I k}{(I + 1)(I + 1 + \delta I)}$.
Furthermore, the variance of $\zeta$ is 
\begin{equation*}
\sigma^2=\mathbb{E}[\zeta^2]-\mu^2=\frac{k\cdot(k+1)}{(I+1)(I+2)}-(\frac{k}{I+1})^2=\frac{k(I-k+1)}{(I+1)^2(I+2)}.
\end{equation*}
According to Bernstein's inequality, for any $p \in (0, 1)$,
$$\Delta\mu \leq \sqrt{\frac{\sigma^2 \log(2/p)}{\theta}} + \frac{2\log(1/p)}{\theta}$$
holds with probability at least $1 - p$. 
Thus, we ensure $\Delta I\leq\delta I$ with probability at least $1 - p$ by setting $\theta$ such that
\begin{equation*}
\sqrt{\frac{\sigma^2 \log(2/p)}{\theta}} + \frac{2\log(1/p)}{\theta} \leq \frac{\delta I k}{(I + 1)(I + 1 + \delta I)}.
\end{equation*}
By solving the above inequality as a quadratic inequality in terms of $\sqrt{\theta}$, $\Delta I \leq \delta I$ holds with probability at least $1-p$ if
\begin{equation*}
\resizebox{\hsize}{!}{
$\theta \geq \bigg( \sqrt{\frac{I-k+1}{I+2} \log{\frac{2}{p}}} + \sqrt{\frac{I-k+1}{I+2} \log{\frac{2}{p}} + 8\frac{\delta I(I+1)}{I+1+\delta I}\log{\frac{1}{p}}}\bigg)^2 \frac{(I+1+\delta I)^2}{4\delta^2 I^2 k}$.
}
\end{equation*}
By simplifying the above inequality, we have $\theta = \Theta(\frac{I}{\delta k}\log{\frac{1}{p}})$.
Since $I \leq n$ for any image, by setting $\theta = \Theta(\frac{n}{\delta k} \log{\frac{1}{p}})$, $\Delta I \leq \delta \cdot I$ holds with probability at least $1 - p$.
\end{proof}
}

\begin{theorem}
\label{theorem:dglobal}
  Let $\theta=\Theta(\frac{\Lambda}{\epsilon k} \log{\frac{|V|}{p}})$ for any $\epsilon\in (0, 1)$, Algorithm~\ref{alg:propagation} correctly answers an $\epsilon$-approximate \ced query under degree centrality with probability at least $1 - p$.
\end{theorem}

\techreport{
\begin{proof}
  Let $\widetilde{N}(u)$ denote the estimated degree of node $u$ via sketch propagation.
  Given any $u\in S^+_0(v^\lambda)$ and $v\in S^-_\epsilon(v^\lambda)$, 
  \begin{align*}
    \widetilde{N}(v) & \leq (1+\delta) \cdot N(v) \leq (1+\delta)(1-\epsilon) \cdot N(v^\lambda) \\
                     & \leq (1+\delta)(1-\epsilon) \cdot N(u) \leq \frac{(1+\delta)(1-\epsilon)}{1-\delta} \cdot \widetilde{N}(u)
  \end{align*}
  hold by applying Lemma~\ref{theorem:dconcentration} over $\kset(\iset(v))$ and the definitions of $S^-_\epsilon(v^\lambda)$ and $S^+_0(v^\lambda)$.
  By setting $\delta<\frac{\epsilon}{2-\epsilon}$, we have $\widetilde{N}(v)<\widetilde{N}(u)$, i.e., the sketches rank $v$ lower than $u$, with probability at least $1 - p$.
  Taking the union bound over all nodes in $S^+_0(v^\lambda)$ and setting $\theta = \Theta(\frac{\Lambda}{\epsilon k}\log{\frac{|V|}{p}})$, the sketches rank all nodes in $S^+_0(v^\lambda)$ higher than $v\in S^-_\epsilon(v^\lambda)$. Since there are at least $\lambda|V|$ nodes in $S^+_0(v^\lambda)$, $v$ will not be included in the result $S$.
  Similarly, given any nodes $u \in S^-_0(v^\lambda)$ and $v \in S^+_\epsilon(v^\lambda)$, we have $\widetilde{N}(v) \geq (1-\delta) \cdot N(v) \geq (1-\delta)(1+\epsilon) \cdot N(v^\lambda) \geq (1-\delta)(1+\epsilon) \cdot N(u) \geq \frac{(1-\delta)(1+\epsilon)}{1+\delta} \cdot \widetilde{N}(u)$.
  By setting $\delta<\frac{\epsilon}{2+\epsilon}$, it holds that $\widetilde{N}(v)>\widetilde{N}(u)$, i.e. the sketches rank $v$ higher than $u$, with probability at least $1 - p$.
  Also taking the union bound over all nodes in $S^-_0(v^\lambda)$ and setting $\theta = \Theta(\frac{\Lambda}{\epsilon k}\log{\frac{|V|}{p}})$, the sketches rank all nodes in $S^-_0(v^\lambda)$ lower than $v \in S^+_\epsilon(v^\lambda)$. Since there are at least $(1-\lambda) |V|$ nodes in $S^-_0(v^\lambda)$, $v$ will be included in $S$.
  Finally, by taking the union bound over all nodes in $S^+_\epsilon(v^\lambda)$ and $S^-_\epsilon(v^\lambda)$ and setting $\theta=\Theta(\frac{\Lambda}{\epsilon k}\log{\frac{|V|^2}{p}}) = \Theta(\frac{\Lambda}{\epsilon k}\log{\frac{|V|}{p}})$, all nodes in $S^+_\epsilon(v^\lambda)$ are included in $S$, whereas all nodes in $S^-_\epsilon(v^\lambda)$ are excluded from $S$, with probability at least $1 - p$.
\end{proof}
}

The time complexity of Algorithm~\ref{alg:propagation} is $O(\frac{\Lambda |\eset^*|}{\epsilon}\log{\frac{|V|}{p}})$.
Its bottleneck lies in the KMV sketch construction process, which propagates $\theta$ sketches of sizes at most $k$ over the matching graph with $|\eset^*|$ edges.
Its space complexity is $O(k\theta|\vset^*| + |\eset^*|)$. Compared to the exact algorithm, the additional space is used to store KMV sketches.

\subsection{Early Termination for Personalized \ced}
\label{sec:degree-personal}

Given a query node $q$, a personalized \ced query can be answered directly using the estimated node degrees through \emph{sketch propagation}.
Specifically, if $q$ is ranked higher than the $(1-\lambda)$-quantile node in terms of estimated degree, it returns \texttt{TRUE} for the personalized query, and \texttt{FALSE} otherwise.
The following corollary is a direct extension of Theorem~\ref{theorem:dglobal} for personalized queries.

\begin{corollary}
  When $\theta = \Theta(\frac{\Lambda}{\epsilon k}\log{\frac{|V|}{p}})$, the sketch propagation correctly answers any personalized $\epsilon$-approximate \ced query under degree centrality with probability at least $1 - p$.
\end{corollary}

Nevertheless, we can optimize personalized \ced query processing by terminating the \emph{sketch propagation} process once we can determine that $q$ is not a hub with high confidence.
First, the following lemma inspires the design of early termination optimization.
\begin{lemma}\label{lemma:earlyterm}
  Given a query node $q$, if there exists a node $u \in \vset^*$ such that $I(\bar{u}) \geq \lambda |V|$ and $I(u) > N(q)$, then the answer to the personalized \ced query is \texttt{FALSE}.
\end{lemma}
\techreport{
\begin{proof}
For any node $u \in \vset^*$, any $v_1 \in \iset(\bar{u})$ and $v_2 \in \iset(u)$ must be neighbors in $H$.
Hence, each node in $\iset(\bar{u})$ has at least $I(u)$ neighbors. Since $I(u) > N(q)$ and $I(\bar{u}) \geq \lambda |V|$, there are at least $\lambda |V|$ nodes with degrees higher than $N(q)$.
\end{proof}
}

Based on Lemma~\ref{lemma:earlyterm}, we can terminate the \emph{sketch propagation} process when the estimations $\widetilde{I}(\bar{u})$ and $\widetilde{I}(u)$ of $I(\bar{u})$ and $I(u)$ for $u \in \vset^*$ by KMV sketches satisfy $\widetilde{I}(\bar{u}) \geq \lambda|V|$ and $\widetilde{I}(u) > N(q)$.
Algorithm~\ref{alg:earlyterm} presents the procedure of \texttt{OptimizedReceive}, an optimized version of \texttt{Receive} at Line~\ref{algline:receive} in Algorithm~\ref{alg:propagation} with early termination.
Specifically, \texttt{OptimizedReceive} returns whether the sketch propagation should end at node $u$.
It calls \texttt{Receive} to construct the KMV sketch for the image of $u$ (Line~\ref{algline:orig-receive}) and estimates $I(\bar{u})$ and $I(u)$ accordingly (Lines~\ref{algline:iestimation-begin}--\ref{algline:iestimation-end}).
Finally, it returns whether the estimated $\widetilde{I}(\bar{u})$ and $\widetilde{I}(u)$ have satisfied the early termination conditions (Lines~\ref{algline:etreturn-begin}--\ref{algline:etreturn-end}).

\begin{algorithm}[t]
\caption{{\ttfamily OptimizedReceive}}\label{alg:earlyterm}
\LinesNumbered
\KwIn{Node $u$, arrays of KMV sketches $\kset$}
\KwOut{If sketch propagation can be terminated}
    \texttt{Receive}$(u, \kset)$\; \label{algline:orig-receive}
    $\widetilde{I} \gets 0$ and $\widetilde{I}^\prime \gets 0$\; \label{algline:iestimation-begin}
    \For{$t=1$ \KwTo $\theta$}{
        \textbf{if} $|\kset[\bar{u}][t]| = k$ \textbf{then} $\widetilde{I} \gets \widetilde{I} + \frac{\max(\kset[\bar{u}][t])}{\theta}$\;
        \textbf{else} $\widetilde{I} \gets \widetilde{I} + \frac{k}{(|\kset[\bar{u}][t]| + 1) \cdot \theta}$\;

        \textbf{if} $|\kset[u][t]| = k$ \textbf{then} $\widetilde{I}^\prime \gets \widetilde{I}^\prime + \frac{\max(\kset[u][t])}{\theta}$\;
        \textbf{else} $\widetilde{I}^\prime \gets \widetilde{I}^\prime + \frac{k}{(|\kset[u][t]| + 1) \cdot \theta}$\;
    }
    $\widetilde{I} \gets k/\widetilde{I} - 1$ and $\widetilde{I}^\prime \gets k/\widetilde{I}^\prime - 1$\;\label{algline:iestimation-end}
    \textbf{if} $\widetilde{I} \geq \lambda|V|$ \textnormal{and} $\widetilde{I}^\prime \geq N(q)$ \textbf{then} \Return{\texttt{TRUE}}\;\label{algline:etreturn-begin}
    \Return{\texttt{FALSE}}\; \label{algline:etreturn-end}
\end{algorithm}

\begin{example}
Continuing with Example~\ref{example:propag}, Fig.~\ref{fig:propagation}(b) demonstrates the sketch propagation process for a personalized \ced query with $q=a_4$ and $\lambda=0.4$. Once we obtain the sketch $\kset(\iset(\bar{p}_1))=\{0.2, 0.4\}$, we have $\widetilde{I}(\bar{p}_1)=\frac{2}{0.4}-1=4$ and $\widetilde{I}(p_1)=\frac{2}{0.6}-1=2.3$. Thus, there are $\widetilde{I}(\bar{p}_1)>\lambda|V|$ nodes with an estimated degree greater than $N(a_4)=2$. Therefore, the sketch propagation process can be terminated early with a decision \texttt{FALSE}.
\end{example}

The following theorem bounds the probability that early termination yields false negative results.
\begin{theorem}\label{theorem:beta:degree}
  Let $\theta = O(\frac{\Lambda}{\epsilon k}\log{\frac{|V|}{p}})$, the probability that a personalized \ced query with answer \texttt{TRUE} is answered with \texttt{FALSE} due to early termination at node $u$ is at most $2\cdot p^{(\beta-1)/\epsilon}$, where $\beta = \min(\frac{\widetilde{I}(\bar{u})}{\lambda|V|}, \frac{\widetilde{I}(u)}{N(q)})$.
\end{theorem}

\techreport{
\begin{proof}
We first observe that for any $\delta > 0$, we have
\begin{equation}\label{equation:BigOtransform}
  O(\frac{\Lambda}{\epsilon k} \log{\frac{|V|}{p}}) = O(\frac{\Lambda}{\delta k}\log{\frac{1}{ p^{\delta/\epsilon}}})
\end{equation}
By Lemma~\ref{theorem:dconcentration} and Eqn.~\ref{equation:BigOtransform}, if $\theta = O(\frac{\Lambda}{\epsilon k} \log{\frac{|V|}{p}})$, we have
\begin{equation}\label{equation:deltarelation}
    \widetilde{I}(\bar{u}) \leq (1+\delta) \cdot I(\bar{u}); \; \widetilde{I}(u) \leq (1 + \delta) \cdot I(u).
\end{equation}
Each holds with probability at least $1- p^{\delta/\epsilon}$. 
According to Eqn.~\ref{equation:deltarelation}, for any $\delta<\beta-1$, we have
\begin{equation}
I(\bar{u})\geq \frac{\widetilde{I}(\bar{u})}{1+\delta}\geq\frac{\beta\lambda|V|}{1+\delta}>\lambda|V|
\end{equation}
and
\begin{equation}
I(u)\geq\frac{\widetilde{I}(u)}{1+\delta}\geq\frac{\beta N(q)}{1+\delta}>N(q)
\end{equation}
holds simultaneously with probability $1- 2\cdot p^{(\beta-1)/\epsilon}$.

Note that the early termination occurring at node $u$ leads to a false negative answer only if $\widetilde{I}(\bar{u}) \geq \lambda|V|$ and $\widetilde{I}(u) > N(q)$, but $I(\bar{u}) < \lambda|V|$ or $I(u)\leq N(q)$. Thus, a personalized \ced query with answer \texttt{TRUE} is incorrectly answered due to an early termination at $u$ with probability at most $2\cdot p^{(\beta-1)/\epsilon}$.
\end{proof}
}